% =====================================================================
%  Preamble for slides.tex
% =====================================================================
%  Sections in this file:
%    1. Fonts and general packages
%    2. Colours and code boxes (minted)
%    3. Annotating code
%    4. Small helpers used by the slides
% =====================================================================

% ---------------------------------------------------------------------
%  1. Fonts and general packages
% ---------------------------------------------------------------------
\usepackage[T1]{fontenc}            % font encoding: T1
\usepackage{notomath}               % Noto Serif for text and maths
\usepackage{amsmath}
\usepackage{inconsolata}            % inconsolata font for code

\usepackage{tikz}
\usetikzlibrary{%
    calc,
    positioning,
    fit,
    arrows,
    arrows.meta,
    tikzmark,                       % \tikzmark, used to annotate code (section 3)
    external,
}
\tikzset{>=stealth}
\usepackage{pgfplots}
\pgfplotsset{compat=1.18}

\usepackage{hologo}                 % \hologo{MiKTeX}
\usepackage[table]{xcolor}
\usepackage{graphicx}               % include graphics
\usepackage[export]{adjustbox}      % adds the `frame' key to \includegraphics
\usepackage{booktabs}               % thick and thin lines in tables
\usepackage{array}
\usepackage{pifont}                 % \ding{51} and \ding{55}: tick and cross
\usepackage{siunitx}                % only for the siunitx example in the closing section
\usepackage{qrcode}                 % QR code generation (also clickable)
\hypersetup{hidelinks}              % beamer already loads hyperref

% ---------------------------------------------------------------------
%  2. Colours and code boxes (minted)
% ---------------------------------------------------------------------
%  Code is highlighted by minted, which hands the code to Pygments (Python).
%  That is why the deck must be compiled with -shell-escape; latexmkrc does
%  this for you, and Overleaf does it by default.
%
%  Compared to listings, minted knows every LaTeX command out of the box,
%  so there is no keyword list to maintain.
%  The price: colours come from a named Pygments style, not from \definecolor.
%  To try another style, change `style=` below.
%  `pygmentize -L styles` lists the available ones.
\usepackage{minted}
\setminted{%
    style=xcode,            % purple commands, green comments
    tabsize=4,
    autogobble,             % strip the indentation the code has in the .tex source
    breaklines,             % wrap long lines instead of running off the box
    numbersep=4mm,          % gap between line numbers and code (see codebox below)
}
% Line numbers: small and grey, like 2025.
\renewcommand{\theFancyVerbLine}{\ttfamily\scriptsize\color{numbercolour}\arabic{FancyVerbLine}}
\definecolor{numbercolour}{HTML}{808080}

% Inline code, coloured like the code boxes: \code|\section{Results}|
% Like all verbatim, it only works in frames marked [fragile],
% and not inside the argument of another command (\only{...}, \note{...}).
\newmintinline[code]{latex}{}

% We use every beamercolor here once, so the colour names below are available.
\usebeamercolor{palette primary}
\usebeamercolor{palette quaternary}
\usebeamercolor{block title}
\usebeamercolor{block body}
\usebeamercolor{block title example}
\usebeamercolor{block body example}
\usebeamercolor{block title alerted}
\usebeamercolor{block body alerted}

\tcbuselibrary{minted}              % tcolorbox's bridge to minted
\tcbuselibrary{raster}              % grids of boxes (the poll slide)

% The four code boxes from 2025, now with minted underneath.
% Their first optional argument is passed to minted, e.g.
%   firstline=3, lastline=7      show only lines 3 to 7
%   highlightlines={5,9}         put a coloured background behind lines 5 and 9
%   escapeinside=||              allow LaTeX between |...|, see section 3
% The second optional argument is passed to tcolorbox (width, height, ...).
\tcbset{%
    % The base of every box in this deck that is not a beamer block.
    % Older versions of beamer's tcolorbox theme (up to 2023) applied their
    % settings to every tcolorbox; newer ones (2026) only to beamer's own blocks.
    % So these boxes set everything they rely on themselves, and look the same
    % with either version:
    deckbox/.style={%
        enhanced,                   % needed for `overlay' (the line-number stripe)
        frame hidden,
        boxrule=0mm,
        arc=1.3mm,
        left=0.05cm, right=0.05cm,  % the padding of beamer's blocks
        top=0.8mm, bottom=0mm,
        toptitle=0.2mm,
        width=\linewidth,           % exactly as wide as the line, no wider
        enlarge left by=0mm,
        enlarge right by=0mm,
    },
    % The light purple background of all code, also usable on its own:
    %   \begin{tcolorbox}[codebg] ... \end{tcolorbox}
    codebg/.style={%
        deckbox,
        colback={block body example.bg},
    },
    codebox/.style={%
        codebg,
        listing engine=minted,
        minted language=latex,
        listing only,
        fonttitle=\normalfont,
        % Spacing, measured from the left edge of the box:
        %   the code starts at 6.5mm (left) + 1mm (tcolorbox's boxsep) = 7.5mm,
        %   the line numbers end numbersep=4mm (in \setminted) before that, at 3.5mm,
        %   and the stripe behind them is 4.5mm wide, so 1mm is left on its right.
        left=6.5mm,
        overlay={%
            \begin{tcbclipinterior}%
                \fill[block title example.bg!20!block body example.bg] (frame.south west)%
                    rectangle ([xshift=4.5mm]frame.north west);%
            \end{tcbclipinterior}%
        },
        % room between a code box and whatever comes after it
        after skip=2.5mm,
    },
}
% Colour behind lines chosen with highlightlines={...}
\colorlet{highlightcolor}{block title example.bg!30!block body example.bg}
\setminted{highlightcolor=highlightcolor}

% Code written directly in the frame (frame needs [fragile]), with a title:
%   \begin{exlisting}[<minted options>]{<title>} ... \end{exlisting}
\newtcblisting{exlisting}[2][]{%
    codebox,
    colbacktitle={block title example.bg},
    coltitle={block title example.fg},
    title={\blockstrut#2},
    minted options={linenos, fontsize=\small, #1},
}

% Same, without a title.
\newtcblisting{exlistingnotitle}[1][]{%
    codebox,
    minted options={linenos, fontsize=\small, #1},
}

% Code read from a file, with a title (default title: Example):
%   \tcbexinputlst[<minted options>][<tcolorbox options>]{<file>}[<title>]
% Reading from a file also works inside \only<...>{...}, so these two are
% what the slides use whenever the code changes between overlays.
\NewTCBInputListing{\tcbexinputlst}{O{} O{} m O{Example}}{%
    codebox,
    listing file={#3},
    colbacktitle={block title example.bg},
    coltitle={block title example.fg},
    title={\blockstrut#4},
    minted options={linenos, fontsize=\small, #1},
    #2
}

% Code read from a file, without a title:
%   \tcbinputlst[<minted options>][<tcolorbox options>]{<file>}
\NewTCBInputListing{\tcbinputlst}{O{} O{} m}{%
    codebox,
    listing file={#3},
    minted options={linenos, fontsize=\small, #1},
    #2
}

% ---------------------------------------------------------------------
%  3. Annotating code
% ---------------------------------------------------------------------
%  The idea: put a \tikzmark BETWEEN two tokens of the code, never around
%  them, so minted still highlights the code normally.
%  Turn on escapeinside=|| for that listing, then write e.g.
%
%      |\tikzmark{a}|\documentclass|\tikzmark{b}|{article}|\tikzmark{c}|
%
%  Each mark is a point on the baseline of the code.
%  \annotatebelow and \annotateabove then draw a bracket between two marks
%  and hang a label off it.
%  Like everything with `remember picture', the drawing needs a second
%  compile run to land in the right place (latexmk does that).
%
%  Use them inside \begin{tikzpicture}[remember picture, overlay].

% How far the bracket sits below the baseline (clear of descenders such as
% p, g, { and [), and above it (clear of capitals and \).
\newlength{\annodepth}\setlength{\annodepth}{1.6mm}
\newlength{\annoheight}\setlength{\annoheight}{5.2mm}
\newlength{\annoarm}\setlength{\annoarm}{1.2mm}     % length of the bracket's arms
\tikzset{annotation/.style={thick, structure.fg}}

% \annotatebelow[<label shift>]{<from mark>}{<to mark>}{<drop>}{<label>}
%   <drop>        how far below the bracket the label hangs
%   <label shift> moves the label sideways (e.g. -1cm), the line bends to follow
\newcommand{\annotatebelow}[5][0pt]{%
    \coordinate (anno-l) at ([xshift=1.5pt, yshift=-\annodepth]pic cs:#2);
    \coordinate (anno-r) at ([xshift=-1.5pt, yshift=-\annodepth]pic cs:#3);
    \coordinate (anno-m) at ($(anno-l)!.5!(anno-r) + (0,-\annoarm)$);
    \draw[annotation] (anno-l) -- ++(0,-\annoarm) -| (anno-r);
    \node[annotation, below, align=center, font=\normalsize] (anno-label)
        at ($(anno-m) + (#1,-#4)$) {#5};
    \draw[annotation] (anno-m) |- ($(anno-m)!.5!(anno-label.north)$) -| (anno-label.north);
}

% \annotateabove[<label shift>]{<from mark>}{<to mark>}{<rise>}{<label>}
%   The same, with the bracket above the code and the label above that.
\newcommand{\annotateabove}[5][0pt]{%
    \coordinate (anno-l) at ([xshift=1.5pt, yshift=\annoheight]pic cs:#2);
    \coordinate (anno-r) at ([xshift=-1.5pt, yshift=\annoheight]pic cs:#3);
    \coordinate (anno-m) at ($(anno-l)!.5!(anno-r) + (0,\annoarm)$);
    \draw[annotation] (anno-l) -- ++(0,\annoarm) -| (anno-r);
    \node[annotation, above, align=center, font=\normalsize] (anno-label)
        at ($(anno-m) + (#1,#4)$) {#5};
    \draw[annotation] (anno-m) |- ($(anno-m)!.5!(anno-label.south)$) -| (anno-label.south);
}

% ---------------------------------------------------------------------
%  4. Small helpers used by the slides
% ---------------------------------------------------------------------

% \blockframe{<number>}{<title>}{<one line>}: the slide that opens a block.
\newcommand{\blockframe}[3]{%
    \begin{frame}
        \centering
        {\Large\color{structure.fg} Block #1\par}
        \medskip
        {\Huge #2\par}
        \bigskip
        {\large\color{black!60} #3\par}
    \end{frame}
}

% ---------------------------------------------------------------------
%  The roadmap: the "Schedule" timeline
% ---------------------------------------------------------------------
%  One line from left to right is the whole session.
%  The solid stretches are the four blocks, the dashed ones the breaks.
%  The short upright mark inside each block is where the hands-on part starts.
%
%  Usage, as the content of a frame:
%      \begin{roadmap}{<state>}
%          \roadmapresults
%      \end{roadmap}
%  <state>  0       intro: each block's name and results appear on its own click
%           1 to 4  that block in focus, the ones before it done
%           5       closing: everything done, the results appear one block per click
%  Leave out \roadmapresults to draw only the line and the block names.

% Names and results of the four blocks: edit them here, once.
% \\ breaks a long name into two lines; the results are list items.
\newcommand{\roadmapname}[1]{%
    \ifcase#1\or
        Skeleton\or
        Text\or
        Maths\or
        Illustrations\\\& References\fi
}
\newcommand{\roadmapitems}[1]{%
    \ifcase#1\or
        \item Title
        \item Sections
        \item Table of contents\or
        \item Paragraphs
        \item Line breaks
        \item Lists
        \item Preamble\or
        \item Inline maths
        \item Display maths
        \item Numbered equations
        \item Multi-line equations\or
        \item Tables
        \item Figures
        \item Cross-references
        \item Citations\fi
}

% Lengths along the line, in units of \roadmapunit (roughly minutes).
% A block is 26 units long, its hands-on part starts at unit 19,
% and a break is 6 units (shorter than to scale, it only needs to be visible).
\newlength{\roadmapunit}\setlength{\roadmapunit}{0.112cm}
\colorlet{roadmap-upcoming}{black!25}

\makeatletter
% Where block #1 starts, starts its hands-on part, ends, and has its middle.
\newcommand{\roadmap@positions}[1]{%
    \pgfmathsetmacro{\roadmap@start}{(#1-1)*32}%
    \pgfmathsetmacro{\roadmap@handson}{\roadmap@start+19}%
    \pgfmathsetmacro{\roadmap@end}{\roadmap@start+26}%
    \pgfmathsetmacro{\roadmap@mid}{\roadmap@start+13}%
}

% How block #1 is drawn in the current state: done, focus or upcoming.
\newcommand{\roadmap@setkind}[1]{%
    \def\roadmap@kind{upcoming}%
    \ifnum#1<\roadmap@state\relax \def\roadmap@kind{done}\fi
    \ifnum#1=\roadmap@state\relax \def\roadmap@kind{focus}\fi
    \ifnum\roadmap@state=0 \def\roadmap@kind{done}\fi
}

% From which overlay on the name of block #1 is shown (\roadmap@from).
% Only the intro reveals the names, one block per click after the bare line.
\newcommand{\roadmap@setfrom}[1]{%
    \def\roadmap@from{1}%
    \ifnum\roadmap@state=0 \pgfmathtruncatemacro{\roadmap@from}{#1+1}\fi
}

\NewEnviron{roadmap}[1]{%
    \def\roadmap@state{#1}%
    \hyphenpenalty=10000 % no hyphenated words in names and results
    \begin{center}
    \begin{tikzpicture}[x=\roadmapunit]
        % The same size in every state, so the line never moves between slides.
        \path[use as bounding box] (0,-2.8) rectangle (122,1.4);
        % the breaks
        \foreach \i in {1,2,3} {
            \draw[roadmap-upcoming, line width=1.2pt, dash pattern=on 2pt off 2pt]
                (\i*32-6+1,0) -- (\i*32-1,0);
        }
        % the blocks
        \foreach \i in {1,2,3,4} {
            \roadmap@positions{\i}
            \roadmap@setkind{\i}
            \roadmap@setfrom{\i}
            \csname roadmap@draw@\roadmap@kind\endcsname{\i}
        }
        \BODY
    \end{tikzpicture}
    \end{center}
}

% A block that is done (and every block on the intro): theme colour.
\newcommand{\roadmap@draw@done}[1]{%
    \draw[structure.fg, line width=2pt] (\roadmap@start,0) -- (\roadmap@end,0);
    \draw[structure.fg, line width=1pt] (\roadmap@handson,-1.5mm) -- (\roadmap@handson,1.5mm);
    \only<\roadmap@from->{%
        \node[above, align=center] at (\roadmap@mid,.25) {\roadmapname{#1}};
    }
}
% The block in focus: a thicker line and a big coloured name.
\newcommand{\roadmap@draw@focus}[1]{%
    \draw[structure.fg, line width=4pt] (\roadmap@start,0) -- (\roadmap@end,0);
    \draw[structure.fg, line width=1.2pt] (\roadmap@handson,-2.2mm) -- (\roadmap@handson,2.2mm);
    \node[above, align=center, font=\Large\bfseries, text=structure.fg]
        at (\roadmap@mid,.3) {\roadmapname{#1}};
}
% A block still to come: grey.
\newcommand{\roadmap@draw@upcoming}[1]{%
    \draw[roadmap-upcoming, line width=2pt] (\roadmap@start,0) -- (\roadmap@end,0);
    \draw[roadmap-upcoming, line width=1pt] (\roadmap@handson,-1.5mm) -- (\roadmap@handson,1.5mm);
    \node[above, align=center, font=\small, text=black!45] at (\roadmap@mid,.25) {\roadmapname{#1}};
}

% The results, as a list under each block, starting where the block starts.
% In focus: only that block's list, at normal size.
% Intro and closing: every block's list, small, appearing one block per click.
\newcommand{\roadmapresults}{%
    \ifnum\roadmap@state=0 \roadmap@alllists{2}\fi
    \ifnum\roadmap@state=5 \roadmap@alllists{1}\fi
    \ifnum\roadmap@state>0 \ifnum\roadmap@state<5
        \roadmap@positions{\roadmap@state}
        \node[below right, inner xsep=0pt] at (\roadmap@start,-.35) {%
            \roadmap@list{4.2cm}{\roadmap@state}};
    \fi\fi
}
% Every block's list; block 1 appears on overlay #1, the next one click later, ...
\newcommand{\roadmap@alllists}[1]{%
    \foreach \i in {1,2,3,4} {
        \roadmap@positions{\i}
        \pgfmathtruncatemacro{\roadmap@from}{\i+#1-1}
        \only<\roadmap@from->{%
            \node[below right, inner xsep=0pt] at (\roadmap@start,-.35) {%
                \small\roadmap@list{3.4cm}{\i}};
        }
    }
}
% \roadmap@list{<width>}{<block>}: the block's items as a list.
% Not an itemize: that would indent the bullets, and here the bullets
% have to sit exactly at the start of the block.
\newcommand{\roadmap@list}[2]{%
    \begin{minipage}[t]{#1}
        \raggedright
        \def\item{%
            \par\noindent\hangindent=1.1em\hangafter=1
            \makebox[1.1em][l]{\usebeamercolor[fg]{itemize item}\usebeamertemplate{itemize item}}%
        }%
        \roadmapitems{#2}\par
    \end{minipage}%
}
\makeatother

% \handson{<block>}{<TODOs>}{<minutes>}{<checkpoint>}{<what to do>}
%   The practice slide at the end of every block.
%   Use it as the whole content of a frame.
\newcommand{\handson}[5]{%
    {\fontsize{40}{40}\selectfont\bfseries #3 min}%
    \hspace{1.5em}%
    {\Large\color{black!60} Block #1: \texttt{TODO} #2}

    \bigskip
    {\large #5\par}

    \vfill
    \begin{alertblock}{Stuck?}
        Copy \texttt{checkpoints/#4} over \texttt{workshop.tex} and continue from there.
    \end{alertblock}
}

% A thin grey frame, for boxes based on deckbox (which hides the frame):
%   \begin{tcolorbox}[deckbox, thinframe] ... \end{tcolorbox}
\tcbset{%
    thinframe/.style={%
        colframe=black!25,
        boxrule=.6pt,
        frame style={fill=black!25},
    },
}

% \outputrender[<includegraphics options>]{<file>}: a pre-compiled output,
% framed like the crops on the Sections slide.
% Every example output on the slides is compiled separately (see renders/),
% so text and maths look exactly like the students' own PDF.
% Snippet renders all have the same width (the article's text width),
% so scaling them to the line width gives the same text size on every slide.
\newcommand{\outputrender}[2][]{%
    \includegraphics[width=\dimexpr\linewidth-0.8pt\relax, frame, #1]{#2}%
}

% A box with a compiled-output look: white, thin frame, serif text.
% Being replaced by \outputrender; remove it once no block uses it anymore.
%   \begin{outputbox} A freshly poured coffee... \end{outputbox}
% Each box starts the equation counter at 0, like a fresh document,
% so its first numbered equation is always (1).
\newtcolorbox{outputbox}[1][]{%
    deckbox,
    thinframe,
    after skip=2.5mm,
    colback=white,
    arc=1mm,
    before upper={\setcounter{equation}{0}},
    #1
}

% A red cross and a green tick, for "do / don't" lists.
\colorlet{goaldone}{green!60!black}
\newcommand{\nope}{\textcolor{psecbg}{\ding{55}}}
\newcommand{\yep}{\textcolor{goaldone}{\ding{51}}}
